% ===================================================================
%  도표 제 7 호 — 겨울의 마을. 봄의 시골집.
%
%  Traduction, faite en 2026, en coreen d'aujourd'hui, sur l'ido de
%  texto/io. Regles de la colonne : voir 10-dopyo-01.tex. Deux scenes,
%  deux numerotations : les cles portent c1 et c2.
%
%  LE MARRON CHAUD (42), ET C'EST L'INVERSE DE TOUT CE QU'ON A VU. Ce
%  renvoi est le trou le plus documente du releve : le telougou n'a
%  pas de mot parce que le chataignier ne pousse pas en Andhra, et
%  trois autres colonnes ont du le contourner de meme. Le coreen a le
%  mot, 밤 ; il a aussi la CHOSE, 군밤, les marrons grilles qu'on
%  vend dans la rue en hiver ; et il a jusqu'a la scene entiere du
%  tableau — l'echoppe, le froid, les enfants qui s'en rechauffent les
%  mains. Rien a declarer, rien a contourner : c'est la premiere fois
%  qu'une colonne rencontre cet objet-la et n'a rien a en dire.
%
%  LA BLAGUE DU VIEUX CORDONNIER, ET LE MEME ARBITRAGE QU'AU TELOUGOU.
%  « En ville j'etais shuifisto et je n'avais pas de clients ; ici je
%  suis shureparisto et le travail ne manque pas. » Le coreen a un mot
%  ancien pour le cordonnier : 갖바치. C'est aussi le nom d'un des
%  metiers 천민, encore employe comme insulte, et la colonne ne
%  l'ecrit pas. Elle ecrit 구두 만드는 사람 et 구두 수선하는 사람 —
%  la fonction, deux fois — et la blague tient toujours, parce
%  qu'elle repose sur le registre et non sur le mot. Meme conclusion
%  qu'au tableau 6 telougou, dans une famille de langues sans rapport
%  avec la sienne.
%
%  LA FERME COREENNE, ELLE, EST ENTIERE EN COREEN : 여물시렁 le
%  ratelier, 깔짚 la litiere, 두레박 le seau du puits, 구유
%  l'auge, 써레 la herse, 곡괭이 la pioche, 물레 le rouet, 실톳대 la
%  quenouille, 두엄 le fumier, 볏 la crete, 망아지 le poulain, 송아지
%  le veau, 병아리 le poussin. Les betes de basse-cour ont toutes leur
%  nom, et le coreen distingue le male de la femelle par un prefixe —
%  수칠면조 / 암오리 — la ou l'ido met -ulo et -ino.
%
%  TRENTE-DEUX PAIRES. Le tableau se partage en deux : la premiere
%  scene raconte des visites et enumere des artisans, la seconde POSE
%  la basse-cour objet par objet. C'est la seconde qui coute.
% ===================================================================

%%K t07-noto-1 noto
\VUnotes{143.2mm}{%
(*) 끼니에 관한 자세한 것은 도표 제 6 호와 제 13 호를 보라.%
}

%%K t07-apar-1 apar
\VUsaut{4.0mm}
\VUcentre{13.2pt}{90}{둘째 계열}
\VUsaut{3.0mm}
\VUfilet{16mm}
\VUsaut{5.6mm}
\VUtitre{15.0pt}{60}{도표 제 7 호}
\VUsaut{3.0mm}

%%K t07-tit-1 sub
\VUcentre{12.0pt}{0}{\textit{첫째 마당.}}
\VUsaut{3.0mm}

%%K t07-tit-2 sub
\VUpk{10.2pt}{40}{겨울의 마을.}
\VUsaut{4.0mm}

%%K t07-c1-01-1 p
1. --- 새해 방학에 나는 아버지를 따라 노부르보에 갔다. 아버지가 장사 일을 몇 가지
보셔야 하는 작은 마을이다.

%%K t07-c1-02-1 p
2. --- 우리는 도시와 그 고을 사이를 오가는 \VUgras{역마차} \textsuperscript{(11)}를
탔다. 그런데 날이 몹시 추워서, 그리 편하지 않은 그 수레의 여행은 별로 즐겁지 않았다.

%%K t07-c1-03-1 p
3. --- 그래서 \VUgras{마부}가 광장에서 수레를 세우는 것을 보았을 때 참으로 기뻤다.
「흰 곰」 \VUgras{여관} \textsuperscript{(2)} 앞이었다. \VUgras{여관 주인}
\textsuperscript{(4)}이 제 집 문지방에서 우리를 기다리며 조용히 \VUgras{담뱃대}를
물고 있었다.

%%K t07-c1-03-2 p
그는 우리에게 들어오라고 했고, 나는 두 번 청하게 하지 않고 화덕에서 타닥거리는 마른
가지 불 앞에 자리를 잡았다. 그때 나는 매서운 \VUgras{북풍}을 실컷 비웃었다. 그 바람은
낡은 \VUgras{간판} \textsuperscript{(3)}을 삐걱거리게 하고, 굴뚝에서 나오는
\VUgras{연기} \textsuperscript{(1)}를 검은 소용돌이로 날려 보냈다.

%%K t07-c1-04-1 p
4. --- 점심때였다. 더 기다리지 않고 우리는 차려 준 소박하지만 맛있는 끼니를 잘 먹었다 (*).

%%K t07-tit-3 sub
\VUpk{10.2pt}{40}{몇 군데 들르다.}
\VUsaut{3.0mm}

%%K t07-c1-05-1 p
5. --- 점심을 먹고 나서 나는 보험 일을 하시는 아버지를 따라 손님들을 찾아갔다. 먼저
\VUgras{대장장이} \textsuperscript{(5)}에게 갔다. 여관 가까이에 사는 사람이다. 그는 말에
편자를 박느라 바빴다. 나는 그 조수의 억센 힘에 놀랐다. 그는 제 가죽 앞치마 위에 말
\VUgras{발굽}을 단단히 붙들고 있었고, 그동안 대장장이는 힘찬 망치질로 \VUgras{편자}
\textsuperscript{(6)}를 박았다.

%%K t07-c1-06-1 p
6. --- 그 다음에는 마을 \VUgras{구두장이} \textsuperscript{(7)} 늙은 야코부스에게 갔다.
그는 훌륭한 \VUgras{장화}를 간판으로 걸어 두었지만, 구멍 난 낡은 구두 한 켤레에 밑창을
대는 일로 만족했다.

%%K t07-c1-06-2 p
그 노인은 웃으며 우리에게 말했다. 「도시에서 나는 \VUgras{구두 만드는 사람}이었고
손님이 없었소. 여기서 나는 \VUgras{구두 수선하는 사람}이고 일이 많다오.」

%%K t07-c1-07-1 p
7. --- 그는 이웃인 \VUgras{이발사} \textsuperscript{(8)} 이야기를 했다. 그 사람의 손님들은
만족하지 않는다. 그의 \VUgras{면도칼}은 늘 이가 빠져 있고 그의 \VUgras{가위}는 언제나
잘 들지 않는다.

그러나 그가 마을에 하나뿐인 이발사라 사람들은 어쩔 수 없이 그에게 수염을 깎고 머리를
자른다. 게다가 그는 인색하고 무뚝뚝한 사람이다.

%%K t07-c1-08-1 p
8. --- 다행히 다른 \VUgras{이웃들}은 그와 같지 않다. 이를테면 \VUgras{수레장이}
\textsuperscript{(9)}는 마음씨 좋고 부지런하고 남을 잘 돕는 사람이다. 그에게 수레를 고치게
하는 것은 즐거운 일이다. 그의 일은 언제나 끝이 야무지다.

%%K t07-c1-09-1 p
9. --- 늙은 야코부스는 \VUgras{마구장이} \textsuperscript{(10)}에 대해서도 불평하지 않았다.
일이 밀릴 때면 그가 이따금 \VUgras{마구} 꿰매는 일을 거들어 준다.

%%K t07-c1-09-2 p
아버지가 그에게 집안 소식을 물었다. 「다들 잘 있습니다. 손녀는 \VUgras{학교}
\textsuperscript{(12)}에 다니지요. 그 아이 \VUgras{여선생님} \textsuperscript{(13)}이 아주
흡족해하십니다. 좋은 \VUgras{학생} \textsuperscript{(14)}이거든요. 못된 제 아들놈과는
다릅니다. 그놈은 놀 생각뿐이라, \VUgras{선생님} \textsuperscript{(15)}이 아무것도 가르치지
못하십니다.」

%%K t07-tit-4 sub
\VUsaut{3.0mm}
\VUpk{10.2pt}{40}{마을 이야기.}
\VUsaut{3.0mm}

%%K t07-c1-10-1 p
10. --- 그 노인은 이어 마을 소식을 들려주었다. 새 학교를 지을 참인데, \VUgras{면사무소}에
들어 있던 학교가 너무 작아졌기 때문이다. 그건 쉬운 일이다. \VUgras{방앗간 주인}의 뜰
한쪽을 사면 되니까. 그 딱한 사람에게는 아주 이로울 것이다. 추위 때문에
\VUgras{물방앗간} \textsuperscript{(16)}의 \VUgras{맷돌}이 더는 밀을 갈지 못한다.
\VUgras{물레바퀴} \textsuperscript{(17)}가 \VUgras{얼음} \textsuperscript{(25)}에 얼어붙어
멈추었기 때문이다. \VUgras{개울} \textsuperscript{(23)}의 얼음이다.

%%K t07-c1-11-1 p
11. --- 「짓는 이야기가 나왔으니, 우리 새 \VUgras{교회} \textsuperscript{(18)}는 보셨소?
눈 이불을 덮고 아주 그림 같지요. \VUgras{까마귀} \textsuperscript{(19)}들은 제 낡은 종탑을
알아보지 못해서, \VUgras{무덤골} \textsuperscript{(20)}의 큰 나무에 쓸쓸히 앉아 있소.」

%%K t07-c1-12-1 p
12. --- 무덤골 이야기에 구두장이는 아버지가 알던 사람들, 지난해에 죽은 사람들을 떠올렸다.
「\VUgras{무덤} \textsuperscript{(21)}이 얼마나 늘었는지요, 선생님. \VUgras{사이프러스}
\textsuperscript{(22)} 아래에 말입니다! 죽음이 우리 늙은 공증인을 베어 갔지요. 온 마을
사람이 그 \VUgras{장례}에 갔습니다.」

%%K t07-c1-13-1 p
13. --- 「저기 개울에서 스케이트를 타는 이가 그 뒤를 이은 사람이오. 방금 나에게
자기 \VUgras{스케이트} \textsuperscript{(26)}의 끊어진 가죽끈을 고치러 왔었소.」

%%K t07-c1-14-1 p
14. --- 「개울 기슭에는 새 \VUgras{들지기} \textsuperscript{(27)}도 있소. 전에 헌병이던
사람인데, 마을 개구쟁이들이 무서워하지요. 그의 \VUgras{각반} \textsuperscript{(29)}과
\VUgras{군모} \textsuperscript{(28)}가 보이기만 하면 아이들은 곧바로 달아난다오.」

%%K t07-c1-15-1 p
15. --- 「아! 저기 늙은 요셉이 빵집에 줄 \VUgras{장작 단 실은 작은 수레}
\textsuperscript{(30)}를 끌고 오는군요. --- 그렇군, 하고 아버지가 말씀하셨다. 저
\VUgras{노새} \textsuperscript{(31)} 멍에를 나도 알아보겠소. 마침 그에게 할 말이 있으니
이 기회를 쓰겠습니다. 안녕히 계십시오, 야코부스 어른.」

%%K t07-c1-16-1 p
16. --- 우리가 막 나오는데, 마을 아이들이 학교에서 나와 \VUgras{얼음지치는 곳}
\textsuperscript{(33)}으로 달려갔다. \VUgras{넘어져} \textsuperscript{(34)} 팔다리가 부러질
위험도 아랑곳없이. 그 가운데 하나가 수레장이에게서 제 \VUgras{썰매} \textsuperscript{(32)}
를 찾아, 거기에 동무 하나를 앉히고 달려서 떠났다.

%%K t07-c1-17-1 p
17. --- 다른 아이들은 \VUgras{눈더미} \textsuperscript{(37)}로 달려들었다. \VUgras{다리}
\textsuperscript{(40)} 가까이에 있는 것이다. 그러고는 \VUgras{눈사람} \textsuperscript{(39)}
을 향해 퍼붓기 시작했다. 어제 세운 눈사람인데, 아이들이 모자와 담뱃대와 어깨에 멘 책가방을
붙여 놓았다.

%%K t07-c1-18-1 p
18. --- 아이들은 살을 에는 바람도 추위도 별로 아랑곳하지 않는다. 대부분은 \VUgras{목도리}
\textsuperscript{(35)}조차 없다. \VUgras{눈뭉치} \textsuperscript{(38)} 싸움을 하고 나서
몸을 녹이는 데는 아주 뜨거운 \VUgras{군밤} \textsuperscript{(42)} 한 줌이면 넉넉하다.
늙은 \VUgras{파는 아낙} 야코비나에게서 산 것인데, 그 여자는 너무 얇은 제
\VUgras{낡은 숄} \textsuperscript{(43)} 아래에서 추위에 떨고 있다.

%%K t07-tit-5 sub
\VUsaut{3.0mm}
\VUcentre{12.0pt}{0}{\textit{둘째 마당.}}
\VUsaut{3.0mm}

%%K t07-tit-6 sub
\VUpk{10.2pt}{40}{봄의 시골집.}
\VUsaut{4.0mm}

%%K t07-c2-01-1 p
1. --- 겨울의 온갖 즐거움에도 우리는 \VUgras{봄}이 돌아오는 것을 깊은 기쁨으로 맞는다.

%%K t07-c2-01-2 p
오월 저녁, 농장 마당은 얼마나 흥겨운 그림인가!

%%K t07-c2-02-1 p
2. --- 지평선에서 \VUgras{해} \textsuperscript{(1)}가 천천히 진다. \VUgras{들}
\textsuperscript{(3)}을 제 붉은 \VUgras{빛살} \textsuperscript{(2)}로 적시면서. 부지런한
\VUgras{밭 가는 사람} \textsuperscript{(6)}은 제 \VUgras{밭} \textsuperscript{(4)}을 서둘러
간다.

%%K t07-c2-02-2 p
그는 제 \VUgras{쟁기} \textsuperscript{(5)}로 — 힘센 \VUgras{말} \textsuperscript{(7)}
을 맨 것이다 — \VUgras{고랑} \textsuperscript{(8)}을 낸다. 그 고랑에 \VUgras{씨 뿌리는 사람}
\textsuperscript{(9)}이 한 줌씩 \VUgras{씨}를 던질 것이고, 그 씨가 금빛 이삭을 낼 것이다.

%%K t07-c2-03-1 p
3. --- \VUgras{풀 베는 이} \textsuperscript{(11)}들은 제 낫을 갖추고 풀밭의 풀을 베었고,
말리는 이들은 \VUgras{마른풀} \textsuperscript{(10)}을 말렸다 — \VUgras{쇠스랑}
\textsuperscript{(12)}과 갈퀴로 뒤집어서. 이제 그들이 돌아온다.

%%K t07-c2-03-2 p
어떤 이들은 소가 끄는 무거운 수레 앞에서 걷는다. 그들은 제 연장을 어깨에 메고 있다. 좀
더 게으른 이들은 향긋한 마른풀 위에 편히 자리를 잡았다.

%%K t07-c2-04-1 p
4. --- 지나가면서 그들은 \VUgras{정원사} \textsuperscript{(16)}에게 다정하게 인사한다.
그는 \VUgras{과수원} \textsuperscript{(13)}에서 \VUgras{과일나무}를 다듬고 \VUgras{배나무}
\textsuperscript{(14)}에 접을 붙이느라 바쁘다. \VUgras{자두나무} \textsuperscript{(15)}는
꽃이 피기 시작했고, 거름을 잘 준 \VUgras{남새밭} \textsuperscript{(17)}에서는 남새와
배추와 상추가 벌써 잘 자라고 있다.

%%K t07-c2-04-2 p
낮은 담 위에서 고운 \VUgras{공작} \textsuperscript{(18)}이 제 꼬리를 편다. 그 훌륭한
깃을 뽐내려는 것이다.

%%K t07-tit-7 sub
\VUpk{10.2pt}{40}{농장 마당.}
\VUsaut{3.0mm}

%%K t07-c2-05-1 p
5. --- \VUgras{마구간} \textsuperscript{(19)} 문이 열려 있어 여물시렁과 \VUgras{깔짚}
\textsuperscript{(23)}이 보인다. \VUgras{암말} \textsuperscript{(21)}의 깔짚인데, 그 말은
\VUgras{마부} \textsuperscript{(20)}가 방금 멍에를 풀어 주었다. 긴 다리가 우스꽝스러운
\VUgras{망아지} \textsuperscript{(22)}가 껑충거리며 제 어미를 따라간다.

%%K t07-c2-06-1 p
6. --- \VUgras{우물} \textsuperscript{(29)} 가까이에서 \VUgras{암소} \textsuperscript{(25)}
들이 물을 마시러 간다 — 나무 \VUgras{구유} \textsuperscript{(26)}에서. 그것이 소들의
물통이다. \VUgras{송아지} \textsuperscript{(24)}는 그 틈을 타 젖을 빨고, \VUgras{황소}
\textsuperscript{(27)}는 여물의 좋은 냄새를 맡으며 즐겁게 울부짖고, 힘센 \VUgras{뿔}
\textsuperscript{(28)}이 난 머리를 자랑스레 쳐든다.

%%K t07-c2-07-1 p
7. --- 모두가 일한다. \VUgras{하녀} \textsuperscript{(32)}는 손으로 우물
\VUgras{두레박줄} \textsuperscript{(31)}을 잡는다. 그 여자는 \VUgras{두레박}
\textsuperscript{(30)}으로 물을 길어 가축에게 먹인다. \VUgras{곡식 광} \textsuperscript{(33)}
가까이에서 한 사람이 짚부스러기를 긁어모으고, \VUgras{살림집} \textsuperscript{(34)} 문
앞에서는 늙은 \VUgras{실 잣는 아낙} \textsuperscript{(35)}이 제 물레와 \VUgras{실톳대}로
옛날처럼 \VUgras{아마}와 \VUgras{삼}을 잣는다.

%%K t07-tit-8 sub
\VUsaut{3.0mm}
\VUpk{10.2pt}{40}{농장 주인.}
\VUsaut{3.0mm}

%%K t07-c2-08-1 p
8. --- \VUgras{농장 주인} \textsuperscript{(36)}이 이 모든 일을 이끌고 일꾼들에게 이르는
말을 한다. 그는 \VUgras{써레} \textsuperscript{(40)}와 \VUgras{곡괭이} \textsuperscript{(39)}를
지붕 밑에 들여놓으라고 이른다. 그것을 \VUgras{개집} \textsuperscript{(38)} 곁에 놓아둔 채
잊었기 때문이다. \VUgras{개} \textsuperscript{(37)}의 집이다.

%%K t07-c2-09-1 p
9. --- 그의 외침에 \VUgras{비둘기} \textsuperscript{(41)}가 놀라 온 날개로
\VUgras{비둘기장} \textsuperscript{(42)}으로 날아들고, \VUgras{암탉} \textsuperscript{(43)}
들도 놀라 \VUgras{닭장} \textsuperscript{(44)}으로 서둘러 들어간다.

%%K t07-c2-10-1 p
10. --- \VUgras{양우리} \textsuperscript{(48)} 앞에 늙은 \VUgras{양치기} \textsuperscript{(45)}
가 있다. 그는 제 \VUgras{양떼} \textsuperscript{(49)}를 몰고 돌아온다. 늘 손에서 놓지 않는

%%K t07-c2-10-1 p suite
\VUgras{지팡이} \textsuperscript{(46)}를 쥐고, 빛바랜 \VUgras{덧옷} \textsuperscript{(47)}
차림으로, 그는 우리로 \VUgras{숫양} \textsuperscript{(50)}, \VUgras{암양}
\textsuperscript{(53)}, \VUgras{새끼 양} \textsuperscript{(54)}, \VUgras{암염소}
\textsuperscript{(51)}, \VUgras{새끼 염소} \textsuperscript{(52)}를 몰아넣는다. 그의 충직한
\VUgras{개} \textsuperscript{(55)} 아르구스가 맨 뒤에 와서, 짖어 대며 뒤처지는 것들을
재촉한다.

%%K t07-c2-11-1 p
11. --- 이 모든 소리와 움직임도 착한 \VUgras{젖소} \textsuperscript{(56)}의 고요함은
어지럽히지 못한다. 그 소는 제 \VUgras{방울} \textsuperscript{(57)}을 딸랑거리고, 그동안
사람들은 \VUgras{외양간} \textsuperscript{(58)} 문 앞에서 젖을 짠다.

%%K t07-c2-12-1 p
12. 그 소는 제 젖을 내주어야 한다는 것을 아는 듯하다. \VUgras{농장 안주인}
\textsuperscript{(60)}이 \VUgras{버터 젓는 통} \textsuperscript{(61)}에서 \VUgras{버터}를
만들 수 있도록, 또 훌륭한 \VUgras{치즈} \textsuperscript{(64)}를 만들 수 있도록 말이다. 그
치즈는 \VUgras{큰 통} \textsuperscript{(63)} 위에 놓인 널판에서 물기를 흘린다.

%%K t07-c2-13-1 p
13. --- \VUgras{젖 다루는 방} \textsuperscript{(59)}은 온통 아주 깨끗하게 지켜진다.
\VUgras{농장 머슴} \textsuperscript{(65)} 덕분인데, 그는 방금 \VUgras{우유 통}
\textsuperscript{(62)}을 씻었다 — 손에 들고 있는 큰 물통에서.

%%K t07-tit-9 sub
\VUsaut{3.0mm}
\VUpk{10.2pt}{40}{닭 치는 곳.}
\VUsaut{3.0mm}

%%K t07-c2-14-1 p
14. --- 두꺼운 나막신을 신은 그는 좀 무겁게 걷는다. 그래서 \VUgras{수칠면조}
\textsuperscript{(68)}와 \VUgras{암오리} \textsuperscript{(66)}와 \VUgras{수오리}
\textsuperscript{(67)}와 \VUgras{거위} \textsuperscript{(69)}를 쫓아 버린다. 거위들은
\VUgras{제 부리} \textsuperscript{(70)}를 크게 벌리고 불안하게 소리친다.

%%K t07-c2-14-2 p
좀 더 싸움을 좋아하는 \VUgras{수탉} \textsuperscript{(71)}은 제 붉은 \VUgras{볏}
\textsuperscript{(72)}을 세우고 제 \VUgras{꼬리} \textsuperscript{(73)} 깃을 털며
「꼬끼오」 하고 우렁차게 운다.

%%K t07-c2-15-1 p
15. --- \VUgras{어미 닭} \textsuperscript{(74)}이 \VUgras{두엄} \textsuperscript{(81)} 위를
헤집는다 — 제 \VUgras{병아리} \textsuperscript{(75)}들과 함께. \VUgras{새끼 돼지}
\textsuperscript{(78)}는 제 어미에게 달아난다. 돼지우리 곁에 보이는 그 커다란
\VUgras{어미 돼지} \textsuperscript{(77)}에게로.

%%K t07-c2-16-1 p
16. --- 제 칸에서 \VUgras{토끼} \textsuperscript{(79)}들이 연한 \VUgras{풀}
\textsuperscript{(80)}을 게걸스레 먹는다. 농장 주인의 딸이 갖다 주는 풀이다. 요아니나는
제 토끼들을 아주 사랑해서, 날마다 닭장에서 갓 낳은 \VUgras{알} \textsuperscript{(76)}을
거두고 나면 제 작은 벗들을 보러 온다.

\VUsaut{6.0mm}
\VUfilet{22mm}
